\documentclass[11pt,a4paper]{article}

% ---- geometry ----
\usepackage[top=2.5cm, bottom=2.5cm, left=2.8cm, right=2.8cm]{geometry}

% ---- fonts and encoding ----
\usepackage[T1]{fontenc}
\usepackage[utf8]{inputenc}
\usepackage{lmodern}

% ---- tables ----
\usepackage{booktabs}
\usepackage{tabularx}
\usepackage{array}
\usepackage{longtable}

% ---- code listings ----
\usepackage{listings}
\usepackage{xcolor}

\definecolor{codebg}{RGB}{248,248,250}
\definecolor{codeframe}{RGB}{210,215,225}
\definecolor{codecomment}{RGB}{100,110,120}
\definecolor{codekw}{RGB}{0,70,180}
\definecolor{codestr}{RGB}{150,40,0}

\lstset{
  backgroundcolor=\color{codebg},
  basicstyle=\small\ttfamily,
  keywordstyle=\color{codekw}\bfseries,
  commentstyle=\color{codecomment}\itshape,
  stringstyle=\color{codestr},
  breaklines=true,
  frame=single,
  rulecolor=\color{codeframe},
  numbers=left,
  numberstyle=\tiny\color{gray},
  numbersep=8pt,
  showstringspaces=false,
  tabsize=4,
  captionpos=b,
  xleftmargin=16pt,
  xrightmargin=4pt,
}

% ---- hyperlinks ----
\usepackage[
  colorlinks=true,
  linkcolor=black,
  urlcolor=blue,
  citecolor=black,
]{hyperref}

% ---- section formatting ----
\usepackage{titlesec}
\titleformat{\section}{\Large\bfseries\sffamily}{}{0em}{}[\titlerule\vspace{2pt}]
\titleformat{\subsection}{\large\bfseries\sffamily}{}{0em}{}
\titleformat{\subsubsection}{\normalsize\bfseries\sffamily}{}{0em}{}

% ---- headers / footers ----
\usepackage{fancyhdr}
\pagestyle{fancy}
\fancyhf{}
\fancyhead[L]{\small\sffamily multi-distro-provisioner}
\fancyhead[R]{\small\sffamily\leftmark}
\fancyfoot[C]{\small\thepage}
\renewcommand{\headrulewidth}{0.4pt}

% ---- spacing ----
\usepackage{parskip}
\setlength{\parskip}{6pt}

% ---- misc ----
\usepackage{enumitem}
\setlist{noitemsep, topsep=4pt}
\usepackage{microtype}

% ---- column type for wide tables ----
\newcolumntype{L}[1]{>{\raggedright\arraybackslash}p{#1}}
\newcolumntype{C}[1]{>{\centering\arraybackslash}p{#1}}

% ================================================================
\begin{document}

% ---- title page ----
\begin{titlepage}
\centering
\vspace*{2.5cm}
{\fontsize{28}{32}\selectfont\bfseries\sffamily multi-distro-provisioner}\\[0.6cm]
{\large\sffamily Technical Reference and Operator Guide}\\[2.2cm]
\rule{0.7\textwidth}{0.4pt}\\[1.2cm]
{\normalsize
A cross-distribution Linux provisioner that configures\\
identical server environments on Ubuntu 22.04 LTS and Fedora 39\\
from a single codebase.
}\\[2cm]
\begin{tabular}{ll}
  \textbf{Project:}  & multi-distro-provisioner \\[4pt]
  \textbf{Targets:}  & Ubuntu 22.04 LTS, Fedora 39 \\[4pt]
  \textbf{Tooling:}  & Vagrant, VirtualBox, Bash, systemd \\[4pt]
\end{tabular}
\vfill
{\small\sffamily\today}
\end{titlepage}

\tableofcontents
\newpage

% ================================================================
\section{Introduction}
\label{sec:intro}
\markboth{Introduction}{}

\subsection{What This Project Does}

\texttt{multi-distro-provisioner} automates the complete configuration of a fresh Linux virtual machine into a production-ready server environment. The same provisioning logic runs on two different Linux distributions — Ubuntu~22.04 and Fedora~39 — from a single repository.

Each provisioned machine receives:

\begin{itemize}
  \item System users with precise sudo and Docker group access (\S\ref{sec:users})
  \item An nginx web server with a virtual host and health-check endpoint (\S\ref{sec:nginx})
  \item A distribution-appropriate firewall (\S\ref{sec:firewall})
  \item Docker Engine with a smoke-test verification (\S\ref{sec:docker})
  \item Custom systemd service units for automated monitoring and scheduled alerting (\S\ref{sec:systemd})
\end{itemize}

\subsection{Why Two Distributions?}

Enterprise Linux environments rarely standardize on a single distribution. Ubuntu (Debian-family) dominates cloud workloads and developer tooling; Fedora and its downstream RHEL/CentOS variants dominate enterprise server, telco, and government deployments. Understanding both is a foundational skill for Linux system administration.

Building a provisioner that handles both exposes every place where the two ecosystems diverge: package management, firewall architecture, default file locations, and service naming conventions. The goal is not to abstract away those differences, but to handle them explicitly and understand why they exist.

% ================================================================
\section{Linux Distributions}
\label{sec:distros}
\markboth{Linux Distributions}{}

\subsection{What a Distribution Is}

The Linux \emph{kernel} manages hardware, memory, processes, and I/O. A \textbf{distribution} (distro) bundles the kernel with a package manager, a set of system utilities, and a release policy. The kernel is identical across distributions; everything built on top of it differs.

\subsection{The Two Major Families}

All mainstream server distributions belong to one of two package families:

\begin{center}
\begin{tabular}{L{3cm}L{2.5cm}L{7cm}}
\toprule
\textbf{Family} & \textbf{Package format} & \textbf{Distributions} \\
\midrule
Debian     & \texttt{.deb} & Debian, Ubuntu, Linux Mint, Pop!\_OS \\
RPM        & \texttt{.rpm} & Fedora, RHEL, CentOS Stream, Rocky Linux, AlmaLinux \\
\bottomrule
\end{tabular}
\end{center}

\subsection{Ubuntu vs.\ Fedora at a Glance}

\begin{center}
\begin{tabular}{L{3.5cm}L{5.5cm}L{5cm}}
\toprule
\textbf{Aspect} & \textbf{Ubuntu 22.04 LTS} & \textbf{Fedora 39} \\
\midrule
Parent project      & Debian                    & Red Hat (upstream of RHEL) \\
Package manager     & \texttt{apt} / \texttt{apt-get} & \texttt{dnf} \\
Package format      & \texttt{.deb}             & \texttt{.rpm} \\
Package database    & \texttt{/var/lib/dpkg}    & \texttt{/var/lib/rpm} \\
Repository config   & \texttt{/etc/apt/sources.list.d/} & \texttt{/etc/yum.repos.d/} \\
Default firewall    & \texttt{ufw}              & \texttt{firewalld} \\
nginx config dir    & \texttt{/etc/nginx/sites-enabled/} + \texttt{conf.d/} & \texttt{/etc/nginx/conf.d/} \\
Service manager     & systemd (same)            & systemd (same) \\
Release cadence     & LTS every 2 years         & Every 6 months \\
Enterprise downstream & Ubuntu itself (Canonical) & RHEL, CentOS Stream \\
\bottomrule
\end{tabular}
\end{center}

% ================================================================
\section{Package Management}
\label{sec:packages}
\markboth{Package Management}{}

\subsection{apt — Advanced Package Tool (Ubuntu/Debian)}

\texttt{apt} resolves and installs \texttt{.deb} packages from repositories listed in \texttt{/etc/apt/sources.list} and files under \texttt{/etc/apt/sources.list.d/}. Adding a third-party repository requires two steps: adding a GPG signing key so apt can verify packages, and adding a source list entry that points to the repository URL.

\begin{lstlisting}[language=bash, caption={Docker repository setup with apt (setup\_docker.sh -- Ubuntu)}]
# Step 1: add Docker's GPG key to the apt keyring
install -m 0755 -d /etc/apt/keyrings
curl -fsSL https://download.docker.com/linux/ubuntu/gpg \
    | gpg --dearmor -o /etc/apt/keyrings/docker.gpg
chmod a+r /etc/apt/keyrings/docker.gpg

# Step 2: add Docker's apt repository
echo \
  "deb [arch=amd64 signed-by=/etc/apt/keyrings/docker.gpg] \
  https://download.docker.com/linux/ubuntu jammy stable" \
  | tee /etc/apt/sources.list.d/docker.list

# Step 3: refresh the local package index, then install
apt-get update -qq
apt-get install -y docker-ce docker-ce-cli containerd.io \
                   docker-buildx-plugin docker-compose-plugin
\end{lstlisting}

\texttt{apt-get update} must run before any install when a new repository has been added, because apt needs to download the package list from that repository before it can find packages in it.

\subsection{dnf — Dandified YUM (Fedora/RHEL)}

\texttt{dnf} replaced \texttt{yum} as the standard package manager for RPM-based systems. It resolves dependencies using a SAT solver, which produces more reliable results than yum's older algorithm. Third-party repositories are added as \texttt{.repo} files in \texttt{/etc/yum.repos.d/}.

\begin{lstlisting}[language=bash, caption={Docker repository setup with dnf (setup\_docker.sh -- Fedora)}]
# Add Docker's repository file
dnf config-manager --add-repo \
    https://download.docker.com/linux/fedora/docker-ce.repo

# Install (dnf fetches the package list automatically -- no separate update needed)
dnf install -y docker-ce docker-ce-cli containerd.io \
               docker-buildx-plugin docker-compose-plugin
\end{lstlisting}

\subsection{Command Comparison}

\begin{center}
\begin{tabular}{L{4.8cm}L{4.5cm}L{4.2cm}}
\toprule
\textbf{Operation} & \textbf{Ubuntu (apt)} & \textbf{Fedora (dnf)} \\
\midrule
Refresh package index          & \texttt{apt-get update}           & not required \\
Install a package              & \texttt{apt-get install -y nginx} & \texttt{dnf install -y nginx} \\
Remove a package               & \texttt{apt-get remove nginx}     & \texttt{dnf remove nginx} \\
Search for a package           & \texttt{apt search nginx}         & \texttt{dnf search nginx} \\
Show package info              & \texttt{apt show nginx}           & \texttt{dnf info nginx} \\
List installed packages        & \texttt{dpkg -l}                  & \texttt{rpm -qa} \\
Which package owns a file      & \texttt{dpkg -S /path/to/file}    & \texttt{rpm -qf /path/to/file} \\
Check if a package is installed & \texttt{dpkg -l nginx}           & \texttt{rpm -q nginx} \\
Upgrade all packages           & \texttt{apt-get upgrade}          & \texttt{dnf update} \\
Add a signing key              & \texttt{gpg -{}-dearmor} to \texttt{/etc/apt/keyrings/} & included in \texttt{.repo} file \\
\bottomrule
\end{tabular}
\end{center}

% ================================================================
\section{Web Server — nginx}
\label{sec:nginx}
\markboth{Web Server --- nginx}{}

\subsection{Installation Differences}

On both distributions, the nginx binary and behaviour are identical. The package name is \texttt{nginx} on both; only the installation command differs (\texttt{apt-get} vs.\ \texttt{dnf}). However, there are two configuration differences between the distributions that require active handling.

\textbf{Ubuntu:} nginx ships with a default virtual host in \texttt{/etc/nginx/sites-enabled/default}. That site listens as \texttt{default\_server} on port 80 and intercepts all requests, returning 404 for any path our config adds (including \texttt{/health}). The provisioner removes it:

\begin{lstlisting}[language=bash]
rm -f /etc/nginx/sites-enabled/default
\end{lstlisting}

\textbf{Fedora:} Fedora's \texttt{nginx.conf} contains a built-in \texttt{server \{ listen 80; server\_name \_; \}} block. nginx logs a warning (\texttt{conflicting server name "\_"}) and silently ignores our \texttt{conf.d/default.conf}. The provisioner replaces \texttt{nginx.conf} entirely with a minimal version that only loads \texttt{conf.d/}:

\begin{lstlisting}[language=bash]
cat > /etc/nginx/nginx.conf << 'EOF'
...
http {
    include /etc/nginx/conf.d/*.conf;
    # no built-in server block
}
EOF
\end{lstlisting}

\subsection{Virtual Host Configuration}

The file \texttt{config/nginx/default.conf} is deployed to \texttt{/etc/nginx/conf.d/default.conf} on both VMs:

\begin{lstlisting}[caption={config/nginx/default.conf}]
server {
    listen 80;
    server_name _;

    root /var/www/html;
    index index.html;

    location / {
        try_files $uri $uri/ =404;
    }

    location /health {
        return 200 "OK\n";
        add_header Content-Type text/plain;
        add_header Access-Control-Allow-Origin *;
    }
}
\end{lstlisting}

The \texttt{/health} endpoint returns HTTP~200 with the body \texttt{OK}. The \texttt{Access-Control-Allow-Origin} header lets the browser dashboard page fetch this endpoint from a different port without triggering a CORS block.

% ================================================================
\section{Firewall Management}
\label{sec:firewall}
\markboth{Firewall Management}{}

\subsection{How Linux Firewalls Work}

All Linux firewalls configure \textbf{netfilter}, the kernel's packet filtering subsystem. The differences are in the management layer on top. From lowest to highest level:

\begin{itemize}
  \item \textbf{iptables} / \textbf{ip6tables} — the original rule-based interface (still present underneath everything)
  \item \textbf{nftables} — the modern replacement for iptables (used internally by firewalld on RHEL~9+)
  \item \textbf{ufw} — a simplified command-line frontend over iptables (Ubuntu)
  \item \textbf{firewalld} — a daemon that manages rules using a zone model (Fedora/RHEL)
\end{itemize}

\subsection{ufw — Uncomplicated Firewall (Ubuntu)}

\texttt{ufw} provides a stateful, default-deny firewall with simple commands. The critical operational rule: \emph{always add the SSH rule before enabling the firewall}, otherwise you will lock yourself out.

\begin{lstlisting}[language=bash, caption={setup\_firewall.sh -- Ubuntu}]
ufw default deny incoming   # reject all inbound traffic by default
ufw default allow outgoing  # allow all outbound

ufw allow ssh               # = ufw allow 22/tcp
ufw allow 80/tcp            # HTTP
ufw allow 443/tcp           # HTTPS

ufw --force enable          # enable (--force skips the interactive prompt)
ufw reload
\end{lstlisting}

Rules added with \texttt{ufw allow} are persistent by default. There is no separate ``save'' step.

\subsection{firewalld — Zone-Based Firewall (Fedora/RHEL)}

\texttt{firewalld} is a daemon that manages rules using a \emph{zone model}. Each network interface belongs to a zone; each zone has a policy. The default zone is \texttt{public}. Rules are defined as named \emph{services} (groups of ports) rather than raw port numbers.

The key flag distinction: \texttt{-{}-permanent} writes the rule to disk but does not apply it immediately; \texttt{-{}-reload} applies what is on disk to the running firewall.

\begin{lstlisting}[language=bash, caption={setup\_firewall.sh -- Fedora}]
systemctl enable --now firewalld

firewall-cmd --add-service=ssh   --permanent
firewall-cmd --add-service=http  --permanent
firewall-cmd --add-service=https --permanent

firewall-cmd --reload    # apply the saved rules to the running firewall
\end{lstlisting}

\subsection{Command Comparison}

\begin{center}
\begin{tabular}{L{4.5cm}L{4.5cm}L{4.5cm}}
\toprule
\textbf{Operation} & \textbf{ufw} & \textbf{firewalld} \\
\midrule
Enable firewall         & \texttt{ufw enable}             & \texttt{systemctl start firewalld} \\
Show current rules      & \texttt{ufw status numbered}    & \texttt{firewall-cmd -{}-list-all} \\
Allow a TCP port        & \texttt{ufw allow 80/tcp}       & \texttt{firewall-cmd -{}-add-port=80/tcp -{}-permanent} \\
Allow by service name   & \texttt{ufw allow http}         & \texttt{firewall-cmd -{}-add-service=http -{}-permanent} \\
Remove a rule           & \texttt{ufw delete allow 80/tcp}& \texttt{firewall-cmd -{}-remove-service=http -{}-permanent} \\
Apply saved rules       & automatic                       & \texttt{firewall-cmd -{}-reload} \\
Check if active         & \texttt{ufw status}             & \texttt{systemctl is-active firewalld} \\
Rules persist after reboot & yes (default)               & only with \texttt{-{}-permanent} \\
Default incoming policy & \texttt{ufw default deny}       & zone-based (public = deny) \\
\bottomrule
\end{tabular}
\end{center}

% ================================================================
\section{Docker}
\label{sec:docker}
\markboth{Docker}{}

Docker Engine is installed from Docker's own repositories (not the distribution's), because the distribution-packaged version is typically several major versions behind. The Docker daemon starts automatically on boot via systemd. Both users (\texttt{devops} and \texttt{appuser}) are added to the \texttt{docker} group so they can run containers without \texttt{sudo}.

Installation is verified with a smoke test:

\begin{lstlisting}[language=bash]
docker run --rm hello-world
\end{lstlisting}

If this command completes without error, the daemon is running, the CLI can talk to it, and image pulling works.

\textbf{Important security note:} the \texttt{docker} group is effectively equivalent to root access, because a group member can mount the host filesystem inside a container. In a real environment, access to the group would be controlled more tightly. For this lab, both users are added for convenience.

% ================================================================
\section{User Management}
\label{sec:users}
\markboth{User Management}{}

\subsection{User Definitions}

Users are defined in \texttt{config/users.conf} in the format \texttt{name:group:shell:sudo}:

\begin{lstlisting}[caption={config/users.conf}]
devops:devops:/bin/bash:yes
appuser:appuser:/bin/bash:no
\end{lstlisting}

The \texttt{create\_users.sh} script reads this file and creates each user. \texttt{devops} gets a \texttt{NOPASSWD} sudoers entry; \texttt{appuser} does not.

\subsection{Why mapfile, Not while read}

\begin{lstlisting}[language=bash]
# Correct -- reads into array, leaves stdin free
mapfile -t users < "${USERS_CONF}"
for line in "${users[@]}"; do ...

# Incorrect -- commands inside the loop (useradd, etc.) consume
# bytes from fd 0, which is connected to the config file
while IFS=: read -r name ...; do
    useradd ...   # useradd consumes from stdin -- corrupts the file read
done < "${USERS_CONF}"
\end{lstlisting}

Commands like \texttt{useradd} may read from standard input when called interactively. Using \texttt{while read} redirects the config file to \texttt{fd~0}, and those commands silently consume lines from it. \texttt{mapfile} loads the entire file into memory first, then closes the file descriptor.

\subsection{Sudoers Configuration}

\begin{lstlisting}[language=bash]
echo "devops ALL=(ALL) NOPASSWD: ALL" \
    > /etc/sudoers.d/devops
chmod 440 /etc/sudoers.d/devops
\end{lstlisting}

Sudoers files must be mode \texttt{440} — world-readable is a security issue, and \texttt{sudo} will refuse to use a file with wrong permissions. Files in \texttt{/etc/sudoers.d/} are loaded by the main \texttt{/etc/sudoers} via \texttt{\#includedir}; editing them directly is safer than modifying the main file (a syntax error in \texttt{/etc/sudoers} locks out all sudo access).

% ================================================================
\section{Service Management — systemd}
\label{sec:systemd}
\markboth{systemd}{}

Unlike package management and firewalls, systemd is \textbf{identical} on Ubuntu and Fedora. Both use the same commands, the same unit file syntax, and the same journal for log collection. This section applies equally to both distributions.

\subsection{Key Concepts}

A \textbf{unit} is a file that describes a resource managed by systemd. Relevant types:

\begin{itemize}
  \item \texttt{.service} — a process (daemon). Can be long-running (\texttt{Type=simple}) or one-shot (\texttt{Type=oneshot}).
  \item \texttt{.timer} — a schedule that activates a service. Replaces cron.
  \item \texttt{.socket} — a network or IPC socket that activates a service on first connection.
\end{itemize}

Unit file locations (in priority order):

\begin{center}
\begin{tabular}{L{5.5cm}L{7.5cm}}
\toprule
\textbf{Path} & \textbf{Purpose} \\
\midrule
\texttt{/etc/systemd/system/}    & Admin-managed. Highest priority. Our files go here. \\
\texttt{/run/systemd/system/}    & Runtime-generated. Discarded on reboot. \\
\texttt{/lib/systemd/system/}    & Package-installed. Lowest priority. \\
\bottomrule
\end{tabular}
\end{center}

A file in \texttt{/etc/systemd/system/} with the same name as a package-installed unit silently overrides it.

\subsection{systemctl Reference}

\begin{center}
\begin{tabular}{L{6.5cm}L{7cm}}
\toprule
\textbf{Command} & \textbf{What it does} \\
\midrule
\texttt{systemctl start nginx}          & Start nginx immediately \\
\texttt{systemctl stop nginx}           & Stop nginx immediately \\
\texttt{systemctl restart nginx}        & Stop, then start \\
\texttt{systemctl reload nginx}         & Reload config without stopping \\
\texttt{systemctl enable nginx}         & Start nginx on every boot \\
\texttt{systemctl disable nginx}        & Remove the boot symlink \\
\texttt{systemctl status nginx}         & Human-readable status and recent log lines \\
\texttt{systemctl is-active nginx}      & Exits 0 if running — scriptable \\
\texttt{systemctl is-enabled nginx}     & Exits 0 if enabled for boot \\
\texttt{systemctl daemon-reload}        & Re-read unit files after editing — always run after copying a unit \\
\texttt{systemctl list-timers}          & Show all active timers and their next trigger time \\
\texttt{systemctl list-units -{}-failed} & Show units that are in a failed state \\
\bottomrule
\end{tabular}
\end{center}

\subsection{journalctl Reference}

\begin{center}
\begin{tabular}{L{6.5cm}L{7cm}}
\toprule
\textbf{Command} & \textbf{What it does} \\
\midrule
\texttt{journalctl -u nginx}                 & All logs for the nginx unit \\
\texttt{journalctl -u nginx -f}              & Follow in real time (like \texttt{tail -f}) \\
\texttt{journalctl -u nginx -n 50}           & Last 50 log lines \\
\texttt{journalctl -u nginx -p err}          & Only error-level entries \\
\texttt{journalctl -u nginx --since "1h ago"} & Logs from the past hour \\
\texttt{journalctl -u nginx --no-pager}      & Print without pagination \\
\texttt{journalctl -b}                       & All logs from the current boot \\
\texttt{journalctl --since today}            & All system logs since midnight \\
\texttt{journalctl -u nginx-monitor -f}      & Follow the monitor service \\
\bottomrule
\end{tabular}
\end{center}

\subsection{Unit Files in This Project}

\subsubsection{nginx-monitor.service}

Runs continuously and checks the \texttt{/health} endpoint every 30 seconds. If the check fails three times in a row, it restarts nginx.

\begin{lstlisting}[caption={systemd/nginx-monitor.service}]
[Unit]
Description=nginx health monitor
After=nginx.service network.target

[Service]
Type=simple
ExecStart=/usr/local/bin/nginx-monitor-loop.sh
Restart=always
RestartSec=10
StandardOutput=journal
StandardError=journal

[Install]
WantedBy=multi-user.target
\end{lstlisting}

\texttt{Restart=always} tells systemd to restart the service if it exits for any reason. \texttt{RestartSec=10} adds a 10-second delay between restarts to prevent rapid restart loops. All output goes to the journal, viewable with \texttt{journalctl -u nginx-monitor}.

\subsubsection{log-alert.timer + log-alert.service}

\begin{lstlisting}[caption={systemd/log-alert.timer}]
[Unit]
Description=Hourly nginx log alert check

[Timer]
OnBootSec=5min
OnUnitActiveSec=1h
Persistent=true

[Install]
WantedBy=timers.target
\end{lstlisting}

\texttt{OnBootSec=5min} — fires 5 minutes after boot (gives services time to start).\\
\texttt{OnUnitActiveSec=1h} — fires every hour after the last activation.\\
\texttt{Persistent=true} — if the machine was off when the timer was scheduled to fire, it fires immediately on the next boot. This is a key advantage over cron: no silent missed runs.

\subsubsection{Why systemd timers instead of cron?}

\begin{center}
\begin{tabular}{L{4.5cm}L{4.5cm}L{4.5cm}}
\toprule
\textbf{Feature} & \textbf{cron} & \textbf{systemd timer} \\
\midrule
Logging              & to syslog (indirect)    & journal (integrated with systemctl) \\
Missed runs          & silently skipped         & \texttt{Persistent=true} runs on next boot \\
Dependencies         & not supported            & \texttt{After=}, \texttt{Wants=} in unit file \\
One-shot vs.\ daemon & separate concepts        & same framework \\
View next run        & not built-in             & \texttt{systemctl list-timers} \\
\bottomrule
\end{tabular}
\end{center}

% ================================================================
\section{Repository Structure}
\label{sec:structure}
\markboth{Repository Structure}{}

\begin{lstlisting}[caption={Top-level directory layout}]
multi-distro-provisioner/
  Vagrantfile            -- defines ubuntu-node and fedora-node VMs
  Makefile               -- shortcuts: make up, make provision-ubuntu, etc.
  config/
    users.conf           -- user definitions (name:group:shell:sudo)
    nginx/default.conf   -- nginx virtual host configuration
  scripts/
    common/              -- distro-agnostic scripts
      utils.sh               -- logging, detect_distro, check_root
      create_users.sh        -- reads users.conf, creates users and groups
      install_nginx.sh       -- installs and configures nginx
      generate_status_page.sh-- generates the per-VM status page at provision time
      setup_docker.sh        -- installs Docker Engine
      setup_systemd.sh       -- copies and enables the service units
      nginx_monitor_loop.sh  -- the monitor loop (runs as a systemd service)
      monitor_logs.sh        -- scans nginx logs, alerts on 5xx errors
      health_check.sh        -- verifies all services are running
      user_report.sh         -- prints user/group/sudo report
    ubuntu/
      provision.sh           -- Ubuntu entry point (called by Vagrant)
      setup_firewall.sh      -- ufw configuration
    fedora/
      provision.sh           -- Fedora entry point
      setup_firewall.sh      -- firewalld configuration
  systemd/
    nginx-monitor.service  -- always-running nginx health monitor
    log-alert.service      -- oneshot log scanner
    log-alert.timer        -- triggers log-alert.service every hour
  web/
    dashboard.html         -- browser dashboard showing both VMs side by side
  docs/
    project-guide.tex      -- this document
\end{lstlisting}

\subsection{Provisioning Phases}

\begin{center}
\begin{tabular}{C{1cm}L{4cm}L{4.5cm}L{4cm}}
\toprule
\textbf{Ph.} & \textbf{Script} & \textbf{Ubuntu} & \textbf{Fedora} \\
\midrule
1 & \texttt{create\_users.sh}    & \texttt{useradd}, \texttt{sudoers.d}  & same \\
2 & \texttt{install\_nginx.sh}   & \texttt{apt-get install nginx}        & \texttt{dnf install nginx} \\
3 & \texttt{setup\_firewall.sh}  & \texttt{ufw allow/enable}             & \texttt{firewall-cmd -{}-add-service} \\
4 & \texttt{setup\_docker.sh}    & Docker apt repository                 & Docker dnf repository \\
5 & \texttt{health\_check.sh}    & bash (identical)                      & same \\
6 & \texttt{monitor\_logs.sh}    & bash (identical)                      & same \\
7 & \texttt{setup\_systemd.sh}   & \texttt{systemctl} (identical)        & same \\
\bottomrule
\end{tabular}
\end{center}

\subsection{Idempotency}

Every script is safe to run multiple times with the same outcome. Patterns used:

\begin{itemize}
  \item \texttt{is\_installed nginx} — skips package installation if already present
  \item \texttt{id "\$username"} — skips \texttt{useradd} if the user exists
  \item \texttt{getent group "\$group"} — skips \texttt{groupadd} if the group exists
  \item \texttt{service\_enabled nginx} — skips \texttt{systemctl enable} if already enabled
\end{itemize}

% ================================================================
\section{Running the Project}
\label{sec:running}
\markboth{Running the Project}{}

\subsection{Prerequisites}

\begin{enumerate}
  \item \textbf{VirtualBox} 7.0 or later — \url{https://www.virtualbox.org/wiki/Downloads}
  \item \textbf{Vagrant} 2.3 or later — \url{https://developer.hashicorp.com/vagrant/downloads}
  \item \textbf{make} — on Windows, install via \texttt{winget install GnuWin32.Make} and add \\ \texttt{C:\textbackslash Program Files (x86)\textbackslash GnuWin32\textbackslash bin} to the system PATH
\end{enumerate}

\subsection{First Run}

\begin{lstlisting}[language=bash, caption={Complete first-run sequence}]
# Clone
git clone https://github.com/EsteeCohen/multi-distro-provisioner
cd multi-distro-provisioner

# Start both VMs (downloads ~1 GB of base images on first run, takes ~5-10 min)
make up

# Provision Ubuntu -- runs all 7 phases
make provision-ubuntu

# Provision Fedora -- same 7 phases, different tools
make provision-fedora
\end{lstlisting}

\subsection{Verifying the Installation}

\begin{lstlisting}[language=bash, caption={Health checks}]
# Ubuntu
make ssh-ubuntu
sudo bash /opt/provisioner/scripts/common/health_check.sh
exit

# Fedora
make ssh-fedora
sudo bash /opt/provisioner/scripts/common/health_check.sh
exit

# Verify nginx from the host (no SSH needed)
curl http://localhost:8080/health   # should print: OK
curl http://localhost:8081/health   # should print: OK
\end{lstlisting}

All seven checks in \texttt{health\_check.sh} should pass:

\begin{itemize}
  \item nginx service is running
  \item nginx config is valid
  \item nginx responds on port 80
  \item Docker daemon is running
  \item docker CLI works
  \item can pull and run a container
  \item firewall is active (ufw or firewalld depending on distro)
\end{itemize}

\subsection{Dashboard}

After provisioning, open a browser on the host machine and navigate to:

\begin{center}
\texttt{http://localhost:8080/dashboard.html}
\end{center}

Or use the Makefile shortcut:

\begin{lstlisting}[language=bash]
make dashboard
\end{lstlisting}

The dashboard shows both VMs side by side and polls \texttt{/health} on each every 30~seconds. The Fedora panel will show as offline if the fedora-node VM is not running.

Each VM also serves its own individual status page at \texttt{http://localhost:8080} and \texttt{http://localhost:8081}.

\subsection{Makefile Reference}

\begin{center}
\begin{tabular}{L{5.5cm}L{8cm}}
\toprule
\textbf{Command} & \textbf{What it does} \\
\midrule
\texttt{make up}               & Start both VMs \\
\texttt{make down}             & Halt both VMs \\
\texttt{make clean}            & Destroy VMs and delete \texttt{.vagrant/} \\
\texttt{make provision-ubuntu} & Run all provisioning phases on ubuntu-node \\
\texttt{make provision-fedora} & Run all provisioning phases on fedora-node \\
\texttt{make provision-all}    & Provision both VMs in sequence \\
\texttt{make ssh-ubuntu}       & Open a shell in ubuntu-node \\
\texttt{make ssh-fedora}       & Open a shell in fedora-node \\
\texttt{make status}           & Show which VMs are running \\
\texttt{make dashboard}        & Print dashboard URL and open it \\
\texttt{make help}             & Print all available commands \\
\bottomrule
\end{tabular}
\end{center}

\subsection{Useful In-VM Commands}

\begin{lstlisting}[language=bash]
# Full health report
sudo bash /opt/provisioner/scripts/common/health_check.sh

# User and group report
sudo bash /opt/provisioner/scripts/common/user_report.sh

# Run the log scanner manually
sudo bash /opt/provisioner/scripts/common/monitor_logs.sh

# Watch the nginx monitor in real time
journalctl -u nginx-monitor -f

# Check timer schedule
systemctl list-timers

# Ubuntu: show firewall state
sudo ufw status numbered

# Fedora: show firewall state
sudo firewall-cmd --list-all

# Check all services at once
systemctl status nginx docker nginx-monitor log-alert.timer
\end{lstlisting}

% ================================================================
\section{Command Quick Reference}
\label{sec:reference}
\markboth{Command Quick Reference}{}

\subsection{Commands That Are Identical on Both Distributions}

\begin{center}
\begin{longtable}{L{5cm}L{8.5cm}}
\toprule
\textbf{Command} & \textbf{Purpose in this project} \\
\midrule
\endhead
\texttt{useradd -m -s /bin/bash}     & Create devops and appuser with home directory \\
\texttt{groupadd}                    & Create matching user groups \\
\texttt{usermod -aG docker}          & Add users to the docker group \\
\texttt{getent group docker}         & Check if docker group exists before adding \\
\texttt{id \$user}                   & Check if a user exists (idempotency guard) \\
\texttt{chmod 440}                   & Set sudoers file to owner-read, group-read only \\
\texttt{nginx -t}                    & Validate nginx configuration before starting \\
\texttt{curl -sf -m5}                & Health-check the nginx \texttt{/health} endpoint \\
\texttt{systemctl enable}            & Enable a service or timer for boot \\
\texttt{systemctl start}             & Start a service or timer immediately \\
\texttt{systemctl restart}           & Restart a running service \\
\texttt{systemctl is-active}         & Scriptable status check (exit 0 = running) \\
\texttt{systemctl daemon-reload}     & Re-read unit files after copying new ones \\
\texttt{journalctl -u nginx-monitor} & Read nginx monitor service logs \\
\texttt{df -h /}                     & Check root disk usage percentage \\
\texttt{docker run -{}-rm hello-world} & Smoke-test Docker installation \\
\texttt{mapfile -t lines < file}     & Read config file into array (avoids stdin conflict) \\
\texttt{set -euo pipefail}           & Strict error handling in every script \\
\texttt{hostname -f}                 & Get the fully qualified domain name of the machine \\
\texttt{grep PRETTY\_NAME /etc/os-release} & Get human-readable OS description \\
\bottomrule
\end{longtable}
\end{center}

\subsection{Commands That Differ Between Distributions}

\begin{center}
\begin{longtable}{L{3.8cm}L{4.6cm}L{4.6cm}}
\toprule
\textbf{Purpose} & \textbf{Ubuntu} & \textbf{Fedora} \\
\midrule
\endhead
Install nginx           & \texttt{apt-get install nginx}    & \texttt{dnf install nginx} \\
Update package index    & \texttt{apt-get update -qq}       & not required \\
Install Docker          & \texttt{apt-get install docker-ce}& \texttt{dnf install docker-ce} \\
Add Docker repo         & \texttt{/etc/apt/sources.list.d/} & \texttt{dnf config-manager -{}-add-repo} \\
Import Docker GPG key   & \texttt{gpg -{}-dearmor} to keyrings/ & bundled in \texttt{.repo} file \\
Enable firewall         & \texttt{ufw -{}-force enable}     & \texttt{systemctl enable firewalld} \\
Allow port 80           & \texttt{ufw allow 80/tcp}         & \texttt{firewall-cmd -{}-add-service=http} \\
Show firewall rules     & \texttt{ufw status numbered}      & \texttt{firewall-cmd -{}-list-all} \\
Remove default site     & \texttt{rm sites-enabled/default} & replace \texttt{nginx.conf} entirely \\
Check if pkg installed  & \texttt{dpkg -l nginx}            & \texttt{rpm -q nginx} \\
\bottomrule
\end{longtable}
\end{center}

\end{document}